\documentclass[11pt,a4paper]{article}

\usepackage[utf8]{inputenc}
\usepackage[T1]{fontenc}
\usepackage{amsmath,amssymb,amsthm}
\usepackage{enumerate}
\usepackage[shortlabels]{enumitem}
\usepackage{hyperref}
%\usepackage{microtype}
\usepackage{doi}

\newtheorem{theorem}{Theorem}[section]
\newtheorem{proposition}[theorem]{Proposition}
\newtheorem{corollary}[theorem]{Corollary}
\newtheorem{lemma}[theorem]{Lemma}
\theoremstyle{definition}
\newtheorem{definition}[theorem]{Definition}
\newtheorem{hypothesis}[theorem]{Hypothesis}
\newtheorem{openproblem}[theorem]{Open Problem}
\theoremstyle{remark}
\newtheorem{remark}[theorem]{Remark}

\newcommand{\bstar}{\beta^*}

\title{Projective Persistence and the Physical Sub-Spectrum:\\
A Dynamical Selection Criterion for the Capacity Exponent\\[0.5ex]
\large O20 --- Spectral Admissibility Sub-programme}
\author{J.~Beau\\[0.5ex]
\small Independent Researcher}
\date{2026}

\begin{document}

  \maketitle

  \begin{abstract}
    The O12--O19 sequence establishes that the physically relevant
    observable on $\mathrm{Heis}_3(\mathbb{Z}/q\mathbb{Z})$ is the
    canonical pair-level quantity
    $\sigma_{\mathrm{pair}}^{\mathrm{can}}(n)$, with capacity exponent
    $\delta_{\mathrm{pair}} \approx 7.44$ consistent with the
    phenomenological target $\delta_{\mathrm{pair}} \in [7.4, 10.6]$.
    However, the mechanism selecting $[7.4, 10.6]$ as the physical
    sub-spectrum has not yet been identified.
    The present paper proposes a \emph{projective persistence criterion}:
    a fibre $\{c, q-c\}$ is physically admissible if and only if the
    canonical observable reaches a Born--Infeld saturation threshold
    $\sigma_{\mathrm{BI}}$ within the effective cascade window
    $[n_0, n_1]$ defined by the pre-saturation regime.
    Under two structural hypotheses --- that
    $\sigma_{\mathrm{pair}}^{\mathrm{can}}(n)$ is a monotone function of
    an effective fibre-level amplitude~[H1], and that this amplitude
    inherits the Born--Infeld saturation bound
    $A_n^{\max} = c_{\mathrm{BI}}/\sqrt{\lambda_n}$~[H2] --- the persistence
    condition is equivalent to
    $\delta_{\mathrm{pair}} \in [\delta_{\min}, \delta_{\max}]$,
    where the endpoints are determined by $\sigma_{\mathrm{BI}}$,
    $n_0$, and $n_1$.
    The lower bound $\delta_{\min} \approx 7.4$ is shown to be
    structurally consistent with the minimal fibre-level exponent
    $2\,\delta_{c,\min}$ inherited from O16--O18; the upper bound
    $\delta_{\max} \approx 10.6$ retains a phenomenological component,
    imported via the O7 structural relation from the lepton mass window,
    until $\sigma_{\mathrm{BI}}$ is derived from first principles.
    The intrinsic derivation of $\sigma_{\mathrm{BI}}$ from the
    Born--Infeld constraint $|\partial_t \chi_v| \leq c_{\mathrm{BI}}$ is
    identified as the central open problem of the programme.

    \medskip
    \noindent\textbf{Keywords.}
    spectral admissibility, non-injective projection, projective
    persistence, Born--Infeld saturation, Weil representation,
    canonical pair observable, capacity exponent, physical sub-spectrum
  \end{abstract}

  \clearpage
  \tableofcontents

\clearpage
  \section{Introduction}
    \label{sec:intro}

    The spectral admissibility sub-programme of the Cosmochrony
    framework~\cite{Cosmochrony} aims to identify the structural origin
    of the observed hierarchy of effective mass scales from the dynamics
    of a non-injective projection
    $\Pi : \chi \to \mathcal{O}$
    acting on a relational substrate~$\chi$.
    In this framework, physical observables do not correspond to
    individual configurations of~$\chi$, but to equivalence classes
    induced by the projection.
    The central problem is therefore not only to measure admissible
    quantities, but to identify the correct \emph{observable class}
    on which physical scaling laws are defined.

    A sequence of works (O12--O19) has progressively refined this
    identification.
    O12 and O13~\cite{Beau2026a16,Beau2026a17} introduced a
    representation-adapted observable on Heisenberg Cayley graphs, based
    on the growth of the Gram--Schmidt span under breadth-first
    exploration, and established a stable power-law decay with exponent
    $\hat{\delta}_{\mathrm{exact}} \approx 3.72$.
    However, O14 and O15~\cite{Beau2026a18,Beau2026a19} demonstrated
    that this block-level observable, while mathematically correct,
    cannot reproduce the phenomenological target range
    $\delta \in [7.4, 10.6]$, and that no reweighting or aggregation
    procedure can bridge this gap.
    The discrepancy was shown to arise from an observable-class mismatch,
    rather than from a limitation of the measurement.

    O16~\cite{Beau2026a20} resolved this mismatch by identifying
    conjugate pairs $\{c, q-c\}$ as the minimal invariant units of the
    projection.
    At this level, a new observable emerges whose scaling exponent
    satisfies the exact doubling relation
    $\delta_{\mathrm{pair}} = 2\,\delta_c$, yielding
    $\delta_{\mathrm{pair}} \approx 7.44$, in agreement with the lower
    boundary of the phenomenological range.
    O17~\cite{Beau2026a21} proved the conjugation identity $\rho_{q-c} = \overline{\rho_c}$
    and $\delta_{q-c} = \delta_c$; O18~\cite{Beau2026a22} establishes the parity invariance of
    the Born--Infeld candidate action~\cite{Beau2026c} as a covariance of the response family and states
    the identification of the pair $\{c, q-c\}$ with a fibre of $\Pi$ as an open problem,
    so the pairing keeps the status of a structurally motivated hypothesis.
    O19~\cite{Beau2026a23} completed this refinement by constructing a
    canonical normalisation that removes residual pipeline-dependent
    amplitude factors, yielding a unique, convention-free
    observable~$\sigma_{\mathrm{pair}}^{\mathrm{can}}(n)$.

    At this stage, the observable is fixed, but a fundamental question
    remains open.
    While the exponent $\delta_{\mathrm{pair}} \approx 7.44$ lies within
    the phenomenological target range, the mechanism selecting the
    physical sub-spectrum $\delta \in [7.4, 10.6]$ has not yet been
    identified.
    In particular, it is not known whether this window can be derived
    from first principles, or whether it must be regarded as an external
    constraint imported from the lepton mass hierarchy via the O7
    structural relation
    $\beta^* = 1/(\delta + \tfrac{1}{2})$~\cite{Beau2026a11}.

    The present paper addresses this question.
    We propose a \emph{projective persistence criterion} for the
    canonical fibre-level observable, expressing physical admissibility
    as the requirement that the observable reaches a saturation
    threshold within a finite cascade window.
    This criterion is formulated conditionally on two hypotheses:
    that $\sigma_{\mathrm{pair}}^{\mathrm{can}}(n)$ is a monotone
    function of an effective fibre-level amplitude~[H1], and that this
    amplitude inherits a saturation bound from the Born--Infeld
    constraint~[H2].
    Under these hypotheses, the admissible range of exponents is shown
    to arise from a purely structural condition on the cascade dynamics.

    The resulting picture separates clearly three layers of the problem.
    The existence of a stable power-law regime and the value
    $\delta_{\mathrm{pair}} \approx 7.44$ are established results of
    O16--O19~\cite{Beau2026a20,Beau2026a21,Beau2026a22,Beau2026a23}.
    The persistence criterion provides a structural mechanism capable
    of selecting a finite admissible interval of exponents.
    The derivation of the saturation threshold itself from the underlying
    Born--Infeld dynamics remains an open problem, identified explicitly
    in Section~\ref{sec:open}.

    The paper is organised as follows.
    Section~\ref{sec:observable-hierarchy} summarises the progressive
    refinement of the observable from block-level to canonical
    fibre-level form (O12--O19).
    Section~\ref{sec:canonical-normalisation} recalls the canonical
    normalisation construction of O19 and fixes the notational
    conventions used throughout.
    Section~\ref{sec:persistence-criterion} introduces the projective
    persistence criterion, states the two hypotheses explicitly, and
    derives the associated admissible window.
    Section~\ref{sec:discussion} discusses the epistemic status of the
    result and identifies the derivation of the intrinsic saturation
    threshold as the central open problem of the programme.


  \section{From Block Observables to the Canonical Fibre-Level Observable}
    \label{sec:observable-hierarchy}

    The O-series has refined the physically relevant observable through a
    sequence of steps, each correcting the observable \emph{class} rather
    than the measurement itself.
    We summarise this hierarchy here; the purpose is to fix vocabulary
    and to identify $\sigma_{\mathrm{pair}}^{\mathrm{can}}(n)$ as the
    unique convention-free input to which the persistence criterion of
    Section~\ref{sec:persistence-criterion} applies.

    \subsection{Block-level observable}
      \label{sec:block-obs}

      Papers O12 and O13~\cite{Beau2026a16,Beau2026a17} introduced the
      Gram--Schmidt span tracker on Heisenberg Cayley graphs
      $\mathrm{Cay}(G_q, S_q)$ with
      $G_q = \mathrm{Heis}_3(\mathbb{Z}/q\mathbb{Z})$ for prime $q$.
      For a block of central character
      $c \in (\mathbb{Z}/q\mathbb{Z})^\times$, the tracker records at
      each BFS shell $S_n$ the cumulative span
      $\Sigma^{(c)}(n) =
      \dim \operatorname{span}\{\rho_c(g)v_0 : g \in B_n\}$,
      where $B_n$ is the ball of radius $n$ and $v_0 \in \mathbb{C}^q$ is
      a seed vector.
      The normalised observable $\sigma_c(n)$, derived from
      $\Sigma^{(c)}(n)$ by the O12/O13 pipeline, decays as a power law
      \[
        \sigma_c(n) \sim C_c\, n^{-\delta_c}
      \]
      in the pre-saturation regime $|B_n| \leq q^2$.
      Fitting this law over a range of primes $q$ yields the universally
      stable value
      $\hat{\delta}_{\mathrm{exact}} \approx 3.72 \pm
      0.06$~\cite{Beau2026a16,Beau2026a17}.

      Empirically, $\sigma_c(n)$ is strictly decreasing over the
      pre-saturation window (up to negligible fluctuations) for all tested
      primes and initialisations; this monotonicity is inherited by the
      pair observable $\sigma_{\mathrm{pair}}(n)$ and by its canonical
      form $\sigma_{\mathrm{pair}}^{\mathrm{can}}(n)$.

      The cascade index $n$ is not itself a spectral parameter.
      However, the observed power-law regime is empirically consistent with
      an effective spectral scale $\lambda_n$ governing the admissibility
      envelope, in the sense that the Born--Infeld bound
      $A_n^{\max} \sim c_{\mathrm{BI}} / \sqrt{\lambda_n}$ can be expressed as a
      function of the cascade rank $n$ (up to monotone reparametrisation);
      see Hypothesis~\ref{hyp:H2} and Remark~\ref{rem:hyp-status}.

      This block-level observable is mathematically well-defined and
      correctly measures the admissible capacity of the individual Weil
      representation $\rho_c$.
      However, O14 and O15~\cite{Beau2026a18,Beau2026a19} established that
      it does not constitute the physically final observable:
      $\hat{\delta}_{\mathrm{exact}} \approx 3.72$ lies structurally below
      the phenomenological target $\delta \in [7.4, 10.6]$, and O15
      Corollary~4.8 proves that no reweighting of block-level data can
      raise it above this value.
      The obstacle is not a measurement error but an observable-class
      mismatch: individual blocks $\rho_c$ are half-fibres of the
      projection $\Pi$, not its minimal invariant units.

    \subsection{Pair-level observable}
      \label{sec:pair-obs}

      O16~\cite{Beau2026a20} identified the resolution.
      The projection $\Pi : \chi \to \mathcal{O}$ is non-injective, in the
      structural setting of~\cite{Beau2026c}, and the parity involution
      $c \leftrightarrow q-c$ is the candidate fibre structure of the pair
      sector, a structurally motivated hypothesis
      (the fibre identification is stated as an open problem in O18~\cite{Beau2026a22}).
      The physically relevant observable is accordingly defined on
      conjugate pairs $\{c, q-c\}$:
      \begin{equation}
        \label{eq:sigma-pair}
        \sigma_{\mathrm{pair}}(n) := \sigma_c(n)\,\sigma_{q-c}(n).
      \end{equation}
      The key structural fact, verified numerically to machine precision
      in O16, is that $\delta_{q-c} = \delta_c$, so that the pair
      exponent satisfies the exact doubling relation
      \begin{equation}
        \label{eq:delta-doubling}
        \delta_{\mathrm{pair}} = 2\,\delta_c \approx 7.44,
      \end{equation}
      consistent with the phenomenological target $[7.4, 10.6]$.
      A six-fold reduction in the local-slope variance relative to the
      single-block observable, and a power-law fit with $R^2 > 0.9999$,
      confirm that a genuine asymptotic scaling emerges only at the pair
      level~\cite{Beau2026a20}.

    \subsection{Fibre-level interpretation}
      \label{sec:fibre-interp}

      O17~\cite{Beau2026a21} provided the structural derivation of the
      pair observable.
      The exact conjugation identity $\rho_{q-c} = \overline{\rho_c}$,
      combined with norm-invariance of the Gram--Schmidt process, implies
      that the raw span dynamics of $\rho_c$ and $\rho_{q-c}$ are
      identical under equal initial support (O17 Theorem~2.1).
      The factor $r(c,q)$ appearing in the amplitude ratio
      $\sigma_{q-c}(n) = r(c,q)\,\sigma_c(n)$ is consequently a
      normalisation artefact of the pipeline, not an arithmetic invariant
      of the representation (O17 Proposition~3.1).
      The unique structurally robust consequence is $\delta_{q-c} =
  \delta_c$, which is independent of pipeline normalisation.

      O18~\cite{Beau2026a22} clarifies the status of the pairing.
      The Born--Infeld action $S_{\mathrm{BI}}[\chi]$ is even in
      $\chi$ (O18 Lemma~2.1), and parity acting simultaneously on
      configuration and probe is an covariance of the Born--Infeld
      response family (O18 Proposition~2.6).
      Evenness does not, however, force fixed-probe indiscernibility of
      $\chi$ and $-\chi$ (O18 Proposition~2.4), so the parity orbit
      $\{\chi, -\chi\}$ is not forced into the fibres of $\Pi$;
      the fibre identification is a typed classification problem
      (O18 Problem~2.8).
      In the Weil realisation of
      $\mathrm{Heis}_3(\mathbb{Z}/q\mathbb{Z})$, the conjugation identity
      $\rho_{q-c} = \overline{\rho_c}$ is an O17 theorem, recalled as O18
      Theorem~3.1, while the identification of $\{c, q-c\}$ as a physical
      fibre of $\Pi$ is an open bridge (O18 Remark~3.2).
      The identification of O16 therefore keeps the status of a
      structurally motivated hypothesis.

      This identification makes explicit that the passage from blocks to
      pairs is not an empirical refinement but a structural consequence
      of the non-injectivity of $\Pi$: fibres, not blocks, are the
      invariant objects under the projection.

      \begin{remark}[Observable was mis-identified, not mis-measured]
        \label{rem:mis-identified}
        Papers O14 and O15 correctly compute the admissible capacity at the
        level of individual spectral blocks.
        The move to the fibre-level observable of O16--O18 is not a
        correction of O14--O15 but a change of observable class: the
        non-injectivity of $\Pi$ acts at the level of fibres $\Pi^{-1}(y)$,
        not at the level of individual blocks.
        O15 closes the block-level programme; O16--O18 open and populate the
        fibre-level programme.
      \end{remark}

    \subsection{Canonical pair observable}
      \label{sec:canonical-obs}

      We use the terms \emph{pair-level observable} and
      \emph{fibre-level observable} interchangeably throughout this paper,
      the reading of conjugate pairs $\{c, q-c\}$ as fibres of $\Pi$ being
      the structurally motivated hypothesis whose identification O18 states
      as an open bridge (O18 Remark~3.2).

      The pair observable $\sigma_{\mathrm{pair}}(n)$
      of~\eqref{eq:sigma-pair} retains a residual amplitude factor
      $r(c,q) \in \{1,2,3,4\}$ that depends on the Gram--Schmidt
      initialisation parameters $(b_1, b_2)$ of the O12/O13 pipeline.
      O19~\cite{Beau2026a23} resolved this by making explicit the implicit
      normalisation factor $D(c, b_1, b_2, n)$ of the pipeline and
      classifying the dependence structure of $r(c,q)$ via two
      independence tests (Test~P2 and Test~P3).
      Three outcomes are possible: $r(c,q)$ is a pure basis artefact
      (Case~A), a discrete instance invariant (Case~B), or a mixed
      quantity admitting a decomposition
      $r = r_{\mathrm{base}} \cdot r_{\mathrm{inst}}$ (Case~C); the
      canonical normalisation appropriate to each case is constructed
      explicitly.

      \begin{definition}[Canonical pair observable, O19 Definition~7.1]
        \label{def:sigma-can}
        The canonical pair observable is
        \[
          \sigma_{\mathrm{pair}}^{\mathrm{can}}(n)
          := \sigma_c^{\mathrm{can}}(n)\,\sigma_{q-c}^{\mathrm{can}}(n),
        \]
        where $\sigma_c^{\mathrm{can}}(n)$ is the span observable computed
        with the canonical normalisation of O19 Section~7.
      \end{definition}

      By O19 Proposition~7.2,
      $\sigma_{\mathrm{pair}}^{\mathrm{can}}(n)$ is independent of the
      choice of initial basis and block instance, up to corrections that
      are $o(n^{-\delta_{\mathrm{pair}}})$ asymptotically.
      By O19 Proposition~9.1, the capacity exponent
      $\delta_{\mathrm{pair}}$ is unaffected by this canonicalisation: it
      is a property of the scaling law, not of the amplitude
      normalisation.

      The distinction between the raw cumulative span $\Sigma^{(c)}(n)$
      and the normalised observable $\sigma_c(n)$ is maintained
      throughout this paper: $\Sigma^{(c)}(n) \in \mathbb{Z}_{\geq 0}$
      is the integer-valued dimension of the spanned subspace;
      $\sigma_c(n)$ is the normalised quantity derived from it;
      $r(c,q)$ originates at the level of $\Sigma^{(c)}$ and propagates
      into $\sigma_c$; the asymptotic exponent $\delta_c$ is defined at
      the level of $\sigma_c(n)$.

      In particular, any observable differing from
      $\sigma_{\mathrm{pair}}^{\mathrm{can}}(n)$ by a non-trivial
      normalisation factor depending on $(b_1,b_2)$ or on the block
      instance fails to define an intrinsic fibre-level quantity.
      Up to asymptotically negligible corrections,
      $\sigma_{\mathrm{pair}}^{\mathrm{can}}(n)$ is therefore the unique
      observable compatible with the equivalence relation induced by the
      projection $\Pi$, in the sense that it is invariant under all
      transformations preserving fibre equivalence classes.

      The canonical observable $\sigma_{\mathrm{pair}}^{\mathrm{can}}(n)$
      is the starting point of the present paper.
      O20 applies to it the projective persistence criterion formulated in
      Section~\ref{sec:persistence-criterion}; the exponent
      $\delta_{\mathrm{pair}} \approx 7.44$ extracted from its scaling
      law is the quantity whose membership in the physical sub-spectrum
      $[7.4, 10.6]$ requires dynamical explanation.


  \section{Canonical Normalisation of the Pair Observable}
    \label{sec:canonical-normalisation}

    This section recalls the essential results of
    O19~\cite{Beau2026a23} that are used in
    Section~\ref{sec:persistence-criterion}.
    The goal is not to reproduce the construction of O19 but to fix the
    notational conventions and to make explicit the level at which the
    canonicalisation operates, so that
    Section~\ref{sec:persistence-criterion} is self-contained with
    respect to the definitions it requires.

    \subsection{Pipeline dependence at the amplitude level}
      \label{sec:pipeline-dependence}

      The pair observable $\sigma_{\mathrm{pair}}(n)$
      of~\eqref{eq:sigma-pair} is constructed from two single-block
      observables $\sigma_c(n)$ and $\sigma_{q-c}(n)$.
      Each of these is derived from the cumulative span $\Sigma^{(c)}(n)$
      by the O12/O13 pipeline, which involves an implicit normalisation
      factor
      \[
        D(c, b_1, b_2, n)
      \]
      depending on the central character $c$, the Gram--Schmidt
      initialisation parameters $(b_1, b_2)$, and the BFS depth $n$,
      but with an asymptotically constant contribution in $n$ in the
      pre-saturation regime.
      As a consequence, the amplitude of $\sigma_{\mathrm{pair}}(n)$
      carries a residual factor $r(c,q) \in \{1,2,3,4\}$ whose value
      depends on the choice of $(b_1, b_2)$ and, in general, on the
      block instance.

      Two structural facts, established in O17 and confirmed in
      O19~\cite{Beau2026a21,Beau2026a23}, bound the scope of this
      dependence precisely.
      First, $r(c,q)$ is a pipeline artefact, not an arithmetic invariant
      of the Weil representation $\rho_c$: it does not reflect any
      intrinsic property of the conjugate pair $\{c, q-c\}$.
      Second, and crucially, the pipeline dependence is \emph{entirely at
  the amplitude level}: the asymptotic exponent
      $\delta_{\mathrm{pair}}$ is unaffected by any change of
      normalisation convention (O19 Proposition~9.1).
      The canonicalisation problem is therefore decoupled from the
      exponent extraction: $\delta_{\mathrm{pair}} \approx 7.44$ is a
      stable result independent of the normalisation, while
      $\sigma_{\mathrm{pair}}^{\mathrm{can}}(n)$ is needed only when the
      amplitude itself enters a criterion --- as it does in the
      persistence condition of
      Section~\ref{sec:persistence-criterion}.

    \subsection{Canonical normalisation}
      \label{sec:canonical-def}

      O19 makes $D(c, b_1, b_2, n)$ explicit and classifies the
      dependence structure of $r(c,q)$ by two independence tests.
      Test~P2 varies the initial basis $(b_1, b_2)$ at fixed block
      instance; Test~P3 varies the block instance at rigidly fixed basis.
      Three cases arise:
      \begin{enumerate}[label=(\Alph*)]
  \item $r(c,q)$ is a pure basis artefact, eliminable by a
  canonical choice of initial frame;
  \item $r(c,q)$ is a discrete instance invariant, requiring a
  canonical quotient;
  \item $r(c,q) = r_{\mathrm{base}} \cdot r_{\mathrm{inst}}$,
  combining both components, each eliminated by the corresponding
  procedure.
      \end{enumerate}
      In each case, O19 Section~7 constructs the canonical normalisation
      explicitly.
      We import its output as a definition.

      \begin{definition}[Canonical pair observable]
        \label{def:sigma-can-recap}
        For $q$ an odd prime and
        $c \in (\mathbb{Z}/q\mathbb{Z})^\times$, the canonical pair
        observable is
        \[
          \sigma_{\mathrm{pair}}^{\mathrm{can}}(n)
          \;:=\; \sigma_c^{\mathrm{can}}(n)\;\sigma_{q-c}^{\mathrm{can}}(n),
        \]
        where $\sigma_c^{\mathrm{can}}(n)$ is the span observable of O19
        Section~7, constructed under the canonical normalisation appropriate
        to whichever of Cases~A, B, or~C holds.
      \end{definition}

    \subsection{Asymptotic invariance of the exponent}
      \label{sec:exponent-invariance}

      The relation between $\sigma_{\mathrm{pair}}^{\mathrm{can}}(n)$ and
      the raw pair observable $\sigma_{\mathrm{pair}}(n)$ is
      multiplicative: the canonical normalisation replaces $r(c,q)$ by a
      factor that is either unity (Case~A after basis change), a constant
      independent of~$n$ (Case~B after quotient), or a product of two
      constants independent of~$n$ (Case~C after both steps).
      In all cases the multiplicative factor is asymptotically constant
      in $n$, so that
      \[
        \log \sigma_{\mathrm{pair}}^{\mathrm{can}}(n)
        = \log \sigma_{\mathrm{pair}}(n) + O(1), \qquad n \to \infty.
      \]
      The following proposition, proved in O19, is the key fact used in
      Section~\ref{sec:persistence-criterion}.

      \begin{proposition}[Exponent invariance under canonicalisation,
        O19 Proposition~9.1]
        \label{prop:exponent-invariance}
        The capacity exponent
        \[
          \delta_{\mathrm{pair}}
          \;=\; -\lim_{n \to \infty}
          \frac{\Delta \log \sigma_{\mathrm{pair}}^{\mathrm{can}}(n)}
          {\Delta \log n}
        \]
        is unchanged under the canonical normalisation of
        Definition~\textup{\ref{def:sigma-can-recap}}.
        In particular, $\delta_{\mathrm{pair}} \approx 7.44$ is a property
        of the scaling law of the pair observable, independent of the
        pipeline normalisation convention.
      \end{proposition}

      \begin{remark}[Scope of O19 relative to O20]
        \label{rem:O19-scope}
        Proposition~\ref{prop:exponent-invariance} confirms that O19
        operates strictly at the amplitude level.
        The canonical observable $\sigma_{\mathrm{pair}}^{\mathrm{can}}(n)$
        is required in the present paper not to change the exponent, but
        to ensure that the persistence threshold $\sigma_{\mathrm{BI}}$ of
        Section~\ref{sec:persistence-criterion} is evaluated on a
        convention-free quantity: a threshold condition applied to a
        pipeline-dependent amplitude would inherit the arbitrary
        factor~$r(c,q)$ and lose physical meaning.
      \end{remark}


  \section{Projective Persistence Criterion}
    \label{sec:persistence-criterion}

    We now formulate the central criterion of this paper.
    The canonical observable $\sigma_{\mathrm{pair}}^{\mathrm{can}}(n)$
    established in O19~\cite{Beau2026a23} is a pipeline-independent
    quantity; the question is what dynamical condition it must satisfy
    for the fibre $\{c, q-c\}$ to be physically admissible.

    \begin{hypothesis}[Monotone amplitude correspondence]
      \label{hyp:H1}
      The canonical pair observable
      $\sigma_{\mathrm{pair}}^{\mathrm{can}}(n)$ is a strictly monotone
      function of an effective fibre-level amplitude
      $A_n^{\mathrm{fib}}$:
      \[
        \sigma_{\mathrm{pair}}^{\mathrm{can}}(n)
        = \Phi\!\left(A_n^{\mathrm{fib}}(n)\right),
        \qquad \Phi \text{ strictly monotone decreasing.}
      \]
    \end{hypothesis}

    \begin{hypothesis}[Capacity amplitude bound]
      \label{hyp:H2}
      The effective fibre-level amplitude $A_n^{\mathrm{fib}}$ inherits
      the capacity envelope
      \[
        A_n^{\mathrm{fib}} \leq A_n^{\max}
        = \frac{c_{\mathrm{BI}}}{\sqrt{\lambda_n}},
      \]
      where $c_{\mathrm{BI}}$ is the bound of the capacity axiom
      [A-cap] of~\cite{Beau2026c} and
      $\lambda_n$ is the effective spectral parameter at cascade rank~$n$.
      Saturation occurs at $A_n^{\mathrm{fib}} = A_n^{\max}$, defining
      an intrinsic threshold
      \[
        \sigma_{\mathrm{BI}}
        \;:=\; \Phi\!\left(A_n^{\max}\right).
      \]
    \end{hypothesis}

    \begin{remark}[Epistemic status of
    \textup{[H\ref{hyp:H1}]} and \textup{[H\ref{hyp:H2}]}]
      \label{rem:hyp-status}
      Hypothesis~\ref{hyp:H1} is not yet derived from the Gram--Schmidt
      span dynamics of O12/O13~\cite{Beau2026a16,Beau2026a17}.
      What is established is that
      $\sigma_{\mathrm{pair}}^{\mathrm{can}}(n)$ measures the rank
      increment of the spanned subspace under fibre-level projection;
      the identification of this increment with a saturating amplitude
      variable requires a further structural step left as an open problem
      (Section~\ref{sec:open}).
      Hypothesis~\ref{hyp:H2} rests on the modal envelope
      $A_n^{\max} = c_{\mathrm{BI}} / \sqrt{\lambda_n}$ used throughout the
      relaxation-cascade line~\cite{Beau2026a5}: for an oscillatory mode of
      frequency $\sqrt{\lambda_n}$, the capacity bound of axiom [A-cap]
      of~\cite{Beau2026c} applied to the modal rate gives
      $A_n \sqrt{\lambda_n} \leq c_{\mathrm{BI}}$.
      The envelope is therefore conditional on [A-cap] and on the existence
      of a second-order oscillatory sector; neither is derived
      in~\cite{Beau2026c}.
      The extension from modal amplitudes to fibre-level composite
      observables is the missing link; $\sigma_{\mathrm{BI}}$ is therefore
      not yet an explicit function of $(c_{\mathrm{BI}}, \lambda_c,
      \lambda_{q-c})$.
    \end{remark}

    \begin{remark}[Motivation for Hypothesis~\ref{hyp:H1}]
      \label{rem:H1-motivation}
      The identification postulated in Hypothesis~\ref{hyp:H1} is not
      arbitrary.
      $\sigma_{\mathrm{pair}}^{\mathrm{can}}(n)$ is the only scalar
      observable canonically associated with fibre-level redundancy growth
      that is invariant under all admissible normalisations of the
      pipeline (O19 Proposition~7.2).
      A saturating amplitude is the only structural quantity exhibiting
      comparable monotone depletion behaviour under the Born--Infeld
      constraint.
      Hypothesis~\ref{hyp:H1} asserts that these two depletion processes
      are identified; deriving this identification from first principles
      is Open Problem~\ref{op:H1-derivation}.
    \end{remark}

    \begin{proposition}[Projective persistence criterion]
      \label{prop:persistence}
      Under Hypotheses~\textup{\ref{hyp:H1}}
      and~\textup{\ref{hyp:H2}}, let
      $\sigma_{\mathrm{pair}}^{\mathrm{can}}(n) \sim
      C n^{-\delta_{\mathrm{pair}}}$ on the pre-saturation window
      $[n_0, n_1]$.
      The fibre $\{c, q-c\}$ is projectively persistent if and only if
      the canonical observable reaches the Born--Infeld threshold
      $\sigma_{\mathrm{BI}}$ within the effective cascade window
      $[n_0, n_1]$ defined by the pre-saturation regime:
      \[
        \exists\; n^* \in [n_0, n_1] :
        \quad
        \sigma_{\mathrm{pair}}^{\mathrm{can}}(n^*) = \sigma_{\mathrm{BI}}.
      \]
      Under the empirically established power-law ansatz, $n^*$ is
      explicit:
      \[
        n^*(\delta_{\mathrm{pair}})
        = \left(\frac{C}{\sigma_{\mathrm{BI}}}\right)^{1/\delta_{\mathrm{pair}}},
      \]
      and the persistence condition is equivalent to
      \[
        \delta_{\mathrm{pair}}
        \;\in\;
        \left[
          \frac{\log(C/\sigma_{\mathrm{BI}})}{\log n_1},\;
          \frac{\log(C/\sigma_{\mathrm{BI}})}{\log n_0}
        \right]
        \;=:\; [\delta_{\min}, \delta_{\max}].
      \]
    \end{proposition}

    \begin{proof}
      Strict monotonicity of $\Phi$ (Hypothesis~\ref{hyp:H1}) implies
      that $\sigma_{\mathrm{pair}}^{\mathrm{can}}$ is strictly decreasing
      in $n$; under the power-law ansatz it is continuous and surjective
      onto $(0, C n_0^{-\delta_{\mathrm{pair}}})$.
      The intermediate value theorem then guarantees existence and
      uniqueness of $n^*$.
      The persistence window $[\delta_{\min}, \delta_{\max}]$ follows by
      inverting $n^*(\delta_{\mathrm{pair}}) \in [n_0, n_1]$.
    \end{proof}

    \begin{corollary}[Consistency with the measured exponent]
      \label{cor:delta-consistency}
      The measured value $\delta_{\mathrm{pair}} \approx 7.44$ from
      O16--O19~\cite{Beau2026a20,Beau2026a21,Beau2026a22,Beau2026a23}
      is consistent with
      Proposition~\ref{prop:persistence} in the following sense.
      \begin{enumerate}[label=(\roman*)]
  \item The lower bound $\delta_{\min} \approx 7.4$ is not an
  independent dynamical input.
  It coincides with $2\,\delta_{c,\min} \approx 7.44$, the
  minimal fibre-level exponent inherited from the block spectrum
  by the conjugate-pair doubling of
  O16~\cite{Beau2026a20}.
  The measured $\delta_{\mathrm{pair}} \approx 7.44$ therefore
  falls at the lower boundary of the persistence window,
  consistently with it being the smallest structurally realised
  fibre exponent.
  \item The upper bound $\delta_{\max} \approx 10.6$ corresponds,
  via the O7 structural relation
  $\beta^* = 1/(\delta +
  \tfrac{1}{2})$~\cite{Beau2026a11}, to the phenomenological
  lower bound $\beta^*_{\min} \approx 0.09$ on the cascade
  exponent~\cite{Beau2026a7}.
  Its derivation from first principles --- i.e.\ as a consequence
  of $\sigma_{\mathrm{BI}}$ and $[n_0, n_1]$ rather than an
  input from lepton mass data --- remains open~\textup{[H2]}.
  The role of $\sigma_{\mathrm{BI}}$ is not to fit the
  phenomenological window, but to define an intrinsic admissibility
  threshold independently of it; the consistency of
  $\delta_{\mathrm{pair}} \approx 7.44$ with the interval
  $[7.4, 10.6]$ is therefore a non-trivial posterior consistency
  check, not a calibration condition.
      \end{enumerate}
    \end{corollary}

    \begin{openproblem}[Intrinsic derivation of $\sigma_{\mathrm{BI}}$]
      \label{op:sigma-BI}
      The central open problem of this paper is to derive
      $\sigma_{\mathrm{BI}}$ directly from the capacity axiom [A-cap]
      $|\partial_t \chi_v| \leq c_{\mathrm{BI}}$~\cite{Beau2026c},
      without reference to the phenomenological window
      $\beta^* \in (0.09, 0.13)$.
      This requires establishing the missing link of
      Hypothesis~\ref{hyp:H1}: the identification of
      $\sigma_{\mathrm{pair}}^{\mathrm{can}}(n)$ as a monotone function of
      a fibre-level saturating amplitude inheriting the bound
      $A_n^{\max} = c_{\mathrm{BI}}/\sqrt{\lambda_n}$.
    \end{openproblem}


  \section{Discussion and Outlook}
    \label{sec:discussion}

    \subsection{What is established}
      \label{sec:established}

      The present paper rests on a foundation of results that are
      independent of the hypotheses introduced in
      Section~\ref{sec:persistence-criterion}.
      We recall them here to separate what O20 inherits from what it
      contributes.

      The existence of a stable power-law regime
      $\sigma_{\mathrm{pair}}^{\mathrm{can}}(n)
      \sim C n^{-\delta_{\mathrm{pair}}}$ on the pre-saturation window
      $[n_0, n_1]$ is established in
      O16--O19~\cite{Beau2026a20,Beau2026a21,Beau2026a22,Beau2026a23}.
      The value $\delta_{\mathrm{pair}} \approx 7.44$, consistent with the
      lower boundary of the phenomenological target $[7.4, 10.6]$, is
      stable across all tested primes and initialisations.
      The necessity of the pair level --- no aggregation of block data can
      recover the pair-level exponent, and blocks $c$ and $q-c$ carry
      identical dynamics --- is the combined result of O15 (aggregation
      no-go) and O17 (conjugation identity, recalled in O18
      Theorem~3.1)~\cite{Beau2026a19,Beau2026a21,Beau2026a22};
      the reading of pairs as fibres of $\Pi$ remains the open bridge of
      O18 Remark~3.2.
      The asymptotic invariance of $\delta_{\mathrm{pair}}$ under all
      pipeline normalisations is O19
      Proposition~9.1~\cite{Beau2026a23}, recalled here as
      Proposition~\ref{prop:exponent-invariance}.

      These results are not conditional.
      They constitute the input to the present paper, not its conclusions.

    \subsection{What is proposed}
      \label{sec:proposed}

      The central contribution of this paper is the projective persistence
      criterion of Proposition~\ref{prop:persistence} and its conditional
      derivation of an admissible exponent window.

      Under Hypotheses~\ref{hyp:H1} and~\ref{hyp:H2}, the persistence
      criterion translates physical admissibility into a
      threshold-crossing condition: a fibre $\{c, q-c\}$ is projectively
      persistent if and only if the canonical observable reaches the
      Born--Infeld saturation level $\sigma_{\mathrm{BI}}$ within the
      effective cascade window $[n_0, n_1]$ defined by the
      pre-saturation regime.
      This condition is not dynamical in the temporal sense: $n$ is a
      cascade rank, not a physical time parameter, and the criterion
      concerns the existence of a crossing rank $n^*$ rather than an
      evolution equation.
      Under the empirically established power-law ansatz, the persistence
      condition is equivalent to
      $\delta_{\mathrm{pair}} \in [\delta_{\min}, \delta_{\max}]$, where
      the endpoints are determined by $\sigma_{\mathrm{BI}}$, $n_0$,
      and $n_1$.

      Corollary~\ref{cor:delta-consistency} shows that this framework is
      consistent with the established results in two distinct ways.
      The lower bound $\delta_{\min} \approx 7.4$ is not an independent
      dynamical input: it coincides with $2\,\delta_{c,\min}$, the
      minimal fibre-level exponent inherited from the block spectrum by
      the conjugate-pair doubling of O16.
      The upper bound $\delta_{\max} \approx 10.6$ corresponds, via the
      O7 structural relation
      $\beta^* = 1/(\delta + \tfrac{1}{2})$~\cite{Beau2026a11}, to the
      phenomenological lower bound $\beta^*_{\min} \approx 0.09$
      established in O3~\cite{Beau2026a7}.
      The measured value $\delta_{\mathrm{pair}} \approx 7.44$ falls at
      the lower boundary of the persistence window, which is the position
      expected from the pair-level construction of
      O16--O18~\cite{Beau2026a20,Beau2026a21,Beau2026a22}.

    \subsection{What remains open}
      \label{sec:open}

      The conditional character of the results of
      Section~\ref{sec:persistence-criterion} is explicit and intentional.
      Three open problems organise the residual programme.

      \begin{openproblem}[Intrinsic derivation of $\sigma_{\mathrm{BI}}$]
        \label{op:sigma-BI-disc}
        The saturation threshold $\sigma_{\mathrm{BI}}$ is introduced in
        Hypothesis~\ref{hyp:H2} as a quantity fixed by the capacity axiom
        [A-cap] $|\partial_t \chi_v| \leq c_{\mathrm{BI}}$~\cite{Beau2026c},
        independently of the phenomenological window $[7.4, 10.6]$.
        Its explicit derivation as a function of $(c_{\mathrm{BI}}, \lambda_c,
        \lambda_{q-c})$ from the saturation envelope
        $A_n^{\max} = c_{\mathrm{BI}} / \sqrt{\lambda_n}$ remains open.
        Until this derivation is available, the upper bound
        $\delta_{\max} \approx 10.6$ retains a phenomenological component,
        imported via the O7 relation from the lepton mass window
        $\beta^* \in (0.09, 0.13)$~\cite{Beau2026a7,Beau2026a11}.
        This is the central open problem of the programme: deriving
        $\sigma_{\mathrm{BI}}$ from $|\partial_t \chi_v| \leq c_{\mathrm{BI}}$
        without reference to the phenomenological window would convert the
        persistence criterion from a conditionally consistent structure into
        a fully derived prediction, and would establish that the window
        $[7.4, 10.6]$ is a consequence of the admissibility structure rather
        than an input calibrated to lepton mass data.
      \end{openproblem}

      \begin{openproblem}[Monotone amplitude correspondence]
        \label{op:H1-derivation}
        Hypothesis~\ref{hyp:H1} asserts that
        $\sigma_{\mathrm{pair}}^{\mathrm{can}}(n)$ is a strictly monotone
        function of an effective fibre-level amplitude
        $A_n^{\mathrm{fib}}$.
        This identification --- between the rank increment of the
        Gram--Schmidt span and a saturating amplitude variable inheriting
        the Born--Infeld bound --- is the missing structural link between
        the spectral pipeline and the admissibility structure.
        Establishing it would promote both hypotheses to derived results and
        render the persistence criterion a theorem rather than a conditional
        proposition.
      \end{openproblem}

      \begin{openproblem}[Large-$q$ behaviour of the canonical window]
        \label{op:large-q}
        The persistence window $[\delta_{\min}, \delta_{\max}]$ is derived
        under the power-law ansatz on the finite window $[n_0, n_1]$,
        which depends on $q$.
        The stability of this window as $q \to \infty$ --- equivalently,
        whether the persistence criterion selects a $q$-independent
        sub-spectrum in the large-prime limit --- is not established here.
        This question is related to Open Problem~8.2 of
        O19~\cite{Beau2026a23}, concerning the large-$q$ stability of the
        canonical normalisation.
      \end{openproblem}

    \subsection{Perspectives}
      \label{sec:perspectives}

      Three directions extend naturally from the present paper.

      \paragraph{Numerical testability.}
        The persistence criterion of Proposition~\ref{prop:persistence} is
        directly numerically testable, independently of the derivation of
        $\sigma_{\mathrm{BI}}$ from first principles.
        For a given prime $q$, the canonical observable
        $\sigma_{\mathrm{pair}}^{\mathrm{can}}(n)$ can be computed from the
        O12/O13 pipeline with the canonical initialisation of
        O19~\cite{Beau2026a23}.
        Measuring the crossing rank $n^*$ across a range of pairs
        $(c, q-c)$ and primes $q$ provides a direct empirical estimate of
        the effective threshold $\sigma_{\mathrm{BI}}$, and allows a direct
        test of whether the persistence window $[\delta_{\min},
    \delta_{\max}]$ is consistent with the phenomenological target
        $[7.4, 10.6]$ without invoking the O7 relation.

      \paragraph{Extension beyond
        $\mathrm{Heis}_3(\mathbb{Z}/q\mathbb{Z})$.}
        The persistence criterion is formulated for the Weil realisation of
        $\mathrm{Heis}_3(\mathbb{Z}/q\mathbb{Z})$, but its logical structure
        does not depend on this specific group.
        The same criterion applies to any setting in which a canonical
        fibre-level observable satisfying a power-law decay can be
        constructed.
        Whether the persistence window is universal across groups, or
        depends on representation-theoretic data specific to
        $\mathrm{Heis}_3$, is an open question with implications for the
        generality of the $\delta_{\mathrm{pair}} \to \beta^*$ programme.

      \paragraph{Connection to the
        $\delta_{\mathrm{pair}} \to \beta^*$ derivation.}
        The present paper prepares, but does not complete, the fibre-level
        derivation of the cascade exponent $\beta^*$.
        O21 will verify that the structural conditions underlying the O7
        growth law~\cite{Beau2026a11} are satisfied at the fibre level by
        $\sigma_{\mathrm{pair}}^{\mathrm{can}}$, and derive the transfer
        $\delta_{\mathrm{pair}} \approx 7.44 \mapsto \beta^* \approx 0.126$
        as a legitimate consequence rather than a postulate.
        The canonical observable and the persistence criterion established
        here are the two structural prerequisites for that derivation.


      \bibliographystyle{plain}
      \bibliography{cosmochrony-bibliography}

\end{document}
